% Part: first-order-logic
% Chapter: sequent-calculus
% Section: provability-quantifiers

\documentclass[../../../include/open-logic-section]{subfiles}

\begin{document}

\olfileid[id]{fol}{seq}{qpr}
\olsection{\usetoken{S}{derivability} dan Kuantor}

\begin{explain}
  Teorema kelengkapan juga memerlukan agar aturan-aturan kalkulus sekuen
  menghasilkan fakta-fakta tentang~$\Proves$ yang dibuktikan dalam
  bagian ini.
\end{explain}

\begin{thm}
\ollabel{thm:strong-generalization} Jika $c$ adalah simbol konstanta yang
tidak muncul dalam $\Gamma$ ataupun $!A(x)$ dan $\Gamma \Proves !A(c)$,
maka $\Gamma \Proves \lforall[x][!A(x)]$.
\end{thm}

\begin{proof}
Misalkan $\pi_0$ adalah $\Log{LK}$-!!{derivation} untuk
$\Gamma_0 \Sequent !A(c)$ bagi suatu $\Gamma_0 \subseteq \Gamma$ yang
berhingga. Dengan menambahkan inferensi $\RightR{\lforall}$, kita
memperoleh !!a{derivation} untuk
$\Gamma_0 \Sequent \lforall[x][!A(x)]$, sebab $c$ tidak muncul dalam
$\Gamma$ ataupun $!A(x)$ sehingga syarat variabel eigen terpenuhi.
\end{proof}

\begin{prop}
\ollabel{prop:provability-quantifiers}
\begin{tagenumerate}{prvEx,prvAll}
\tagitem{prvEx}{$!A(t) \Proves \lexists[x][!A(x)]$.}{}
\tagitem{prvAll}{$\lforall[x][!A(x)] \Proves !A(t)$.}{}
\end{tagenumerate}
\end{prop}

\begin{proof}
\begin{tagenumerate}{prvEx,prvAll}
\tagitem{prvEx}{Sekuen $!A(t) \Sequent \lexists[x][!A(x)]$
  !!{derivable}:
  \begin{prooftree}
    \Axiom$!A(t) \fCenter !A(t)$
    \RightLabel{\RightR{\lexists}}
    \UnaryInf$!A(t) \fCenter \lexists[x][!A(x)]$
    \end{prooftree}}{}
\tagitem{prvAll}{Sekuen $\lforall[x][!A(x)] \Sequent !A(t)$
  !!{derivable}:
  \begin{prooftree}
    \Axiom$!A(t) \fCenter !A(t)$
    \RightLabel{\LeftR{\lforall}}
    \UnaryInf$\lforall[x][!A(x)] \fCenter !A(t)$
    \end{prooftree}}{}
\end{tagenumerate}
\end{proof}

\end{document}
